% Part: computability
% Chapter: recursive-functions
% Section: pr-functions

\documentclass[../../../include/open-logic-section]{subfiles}

\begin{document}

\olfileid{cmp}{rec}{prf}
\olsection{आदिम पुनरावर्ती फलने}

आदिम पुनरावर्तन आणि संयोजन वापरून आधीच्या फलनांपासून नवीन फलने कशी
परिभाषित करता येतात, ते पुन्हा नोंदवू.

\begin{defn}
  \ollabel{defn:primitive-recursion} $f$ हे $k$-स्थानी फलन ($k\ge 1$)
  आणि $g$ हे $(k+2)$-स्थानी फलन आहे असे समजा. \emph{$f$ आणि $g$ पासून
  आदिम पुनरावर्तनाने परिभाषित केलेले फलन} म्हणजे पुढील समीकरणांनी परिभाषित
  केलेले $(k+1)$-स्थानी फलन~$h$:
  \begin{align*}
    h(x_0,\dots,x_{k-1},0) & =  f(x_0,\dots,x_{k-1}) \\
    h(x_0,\dots,x_{k-1},y+1) & = g(x_0,\dots,x_{k-1}, y, h(x_0,\dots,x_{k-1}, y))
  \end{align*}
\end{defn}

\begin{defn}
  \ollabel{defn:composition} $f$ हे $k$-स्थानी फलन आणि $g_0$, \dots,
  $g_{k-1}$ ही सर्व $n$-स्थानी असलेली $k$ फलने आहेत असे समजा.
  \emph{$f$ चे $g_0$, \dots,~$g_{k-1}$ यांच्याशी संयोजन करून परिभाषित
    केलेले फलन} म्हणजे पुढील समीकरणाने परिभाषित केलेले $n$-स्थानी
  फलन~$h$:
  \[
  h(x_0, \dots, x_{n-1}) =
  f(g_0(x_0, \dots, x_{n-1}), \dots, g_{k-1}(x_0, \dots, x_{n-1})).
  \]
\end{defn}

$\Succ$ आणि, प्रत्येक नैसर्गिक संख्या $n$ व प्रत्येक $i < n$ साठी,
प्रक्षेपण फलने
\[
\Proj{n}{i}(x_0,\dots,x_{n-1}) = x_i,
\]
यांच्याबरोबरच फलन~$\Zero(x) = 0$ चा आपण आदिम पुनरावर्ती फलनांत
समावेश करू.

\begin{defn}
  आदिम पुनरावर्ती फलनांचा संच म्हणजे $\Nat^n$ पासून $\Nat$ कडे जाणाऱ्या
  आणि पुढील कलमांनी विगमनाधारित रीतीने परिभाषित केलेल्या फलनांचा संच:
  \begin{enumerate}
  \item $\Zero$ हे आदिम पुनरावर्ती आहे.
  \item $\Succ$ हे आदिम पुनरावर्ती आहे.
  \item प्रत्येक प्रक्षेपण फलन $\Proj{n}{i}$ हे आदिम पुनरावर्ती आहे.
  \item $f$ हे $k$-स्थानी आदिम पुनरावर्ती फलन आणि $g_0$,
    \dots,~$g_{k-1}$ ही $n$-स्थानी आदिम पुनरावर्ती फलने असतील, तर
    $f$ चे $g_0$, \dots,~$g_{k-1}$ यांच्याशी केलेले संयोजन आदिम
    पुनरावर्ती असते.
  \item $f$ हे $k$-स्थानी आदिम पुनरावर्ती फलन आणि $g$ हे $k+2$-स्थानी
    आदिम पुनरावर्ती फलन असेल, तर $f$ आणि $g$ पासून आदिम पुनरावर्तनाने
    परिभाषित केलेले फलन आदिम पुनरावर्ती असते.
\end{enumerate}
\end{defn}

\begin{explain}
अधिक संक्षेपाने, आदिम पुनरावर्ती फलनांचा संच हा $\Zero$, $\Succ$ आणि
प्रक्षेपण फलने~$\Proj{n}{j}$ यांचा समावेश असलेला आणि संयोजन व आदिम
पुनरावर्तन यांखाली बंद असलेला लघुतम संच आहे.

आदिम पुनरावर्ती फलनांच्या संचाचे वर्णन करण्याचा आणखी एक मार्ग म्हणजे
त्याची ``टप्प्यांच्या'' रूपात व्याख्या करणे. आरंभीच्या फलनांचा---$\Zero$,
$\Succ$ आणि प्रक्षेपणांचा---संच $S_0$ ने दर्शवू. ही टप्पा~$0$ वरील
आदिम पुनरावर्ती फलने आहेत. टप्पा $S_i$ परिभाषित झाल्यावर, आधीच~$S_i$
मध्ये असलेल्या फलनांना संयोजन किंवा आदिम पुनरावर्तन यांपैकी एका क्रियेचा
एकदा उपयोग करून मिळणाऱ्या सर्व फलनांचा संच $S_{i+1}$ असू द्या. मग
\[
S = \bigcup_{i \in \Nat} S_i
\]
हा सर्व आदिम पुनरावर्ती फलनांचा संच आहे.
\end{explain}

$\Add$ हे आदिम पुनरावर्ती फलन आहे, याची खात्री करू.

\begin{prop}
  बेरीज फलन $\Add(x,y) = x+y$ हे आदिम पुनरावर्ती आहे.
\end{prop}

\begin{proof}
$\Add$ ची $f$ आणि~$g$ या दोन फलनांच्या साहाय्याने केलेली व
\olref{defn:primitive-recursion} च्या आकृतिबंधाशी जुळणारी आदिम
पुनरावर्ती व्याख्या आपल्याकडे आधीच आहे:
\begin{align*}
  \Add(x_0, 0) & = f(x_0) = x_0\\
  \Add(x_0, y+1) & = g(x_0, y, \Add(x_0, y)) =
  \Succ(\Add(x_0, y))
\end{align*}
म्हणून $f$ आणि $g$ ही आदिम पुनरावर्ती असतील, तर $\Add$ देखील आदिम
पुनरावर्ती असते. $f(x_0) = x_0 = \Proj{1}{0}(x_0)$, आणि प्रक्षेपण
फलने आदिम पुनरावर्ती मानली जातात; म्हणून $f$ हे आदिम पुनरावर्ती आहे.
फलन $g$ हे पुढीलप्रमाणे परिभाषित केलेले तीन-स्थानी फलन $g(x_0, y, z)$ आहे:
\[
g(x_0, y, z) = \Succ(z).
\]
यावरून $g$ आदिम पुनरावर्ती आहे असे अजून म्हणता येत नाही, कारण $g$ आणि
$\Succ$ ही अगदी समान फलने नाहीत: $\Succ$ हे एकस्थानी आहे, पण $g$ तीन-स्थानी
असणे आवश्यक आहे. मात्र संयोजनाने आपण $g$ ची ``औपचारिक'' व्याख्या अशी
देऊ शकतो:
\[
g(x_0, y, z) = \Succ(\Proj{3}{2}(x_0, y, z))
\]
$\Succ$ आणि $\Proj{3}{2}$ ही आदिम पुनरावर्ती फलने मानली जात असल्यामुळे,
आदिम पुनरावर्ती फलनांपासून संयोजनाने परिभाषित करता येणारे $g$ देखील आदिम
पुनरावर्ती असते.
\end{proof}

\begin{prop}
  \ollabel{prop:mult-pr}
  गुणाकार फलन $\Mult(x,y) = x \cdot y$ हे आदिम पुनरावर्ती आहे.
\end{prop}

\begin{proof}
  सराव.
\end{proof}

\begin{prob}
  $\Mult$ ची आदिम पुनरावर्ती व्याख्या
  \olref[cmp][rec][prf]{defn:primitive-recursion} ने आवश्यक केलेल्या
  रूपात मांडता येते हे दाखवून आणि संबंधित $f$ व~$g$ ही फलने आदिम
  पुनरावर्ती आहेत हे दाखवून \olref[cmp][rec][prf]{prop:mult-pr} सिद्ध करा.
\end{prob}

\begin{ex}
आदिम पुनरावर्ती व्याख्येचे आपले अगदी पहिले उदाहरण पुढीलप्रमाणे होते:
\begin{align*}
h(0) & =  1 \\
h(y+1) & =  2 \cdot h(y).
\end{align*}
हे फलन \olref{defn:primitive-recursion} ने आवश्यक केलेल्या रूपात बसू शकत
नाही, कारण $k=0$. व्याख्येत $1$ आणि $2$ हे स्थिरही येतात. पहिली अडचण
टाळण्यासाठी एक नाममात्र आदान आणू आणि फलन~$h'$ परिभाषित करू:
\begin{align*}
h'(x_0, 0) & =  f(x_0) = 1 \\
h'(x_0, y+1) & =  g(x_0, y, h'(x_0, y)) = 2 \cdot h'(x_0, y).
\end{align*}
फलन $f(x_0) = 1$ हे $\Succ$ आणि $\Zero$ यांपासून संयोजनाने परिभाषित
करता येते: $f(x_0) = \Succ(\Zero(x_0))$. फलन $g$ हे
$g'(z) = 2 \cdot z$ आणि प्रक्षेपणे यांपासून संयोजनाने परिभाषित करता येते:
\begin{align*}
g(x_0, y, z) & = g'(\Proj{3}{2}(x_0, y, z))
\intertext{आणि पुढे $g'$ ची स्वतःची व्याख्या संयोजनाने अशी देता येते:}
g'(z) & = \Mult(g''(z), \Proj{1}{0}(z))
\intertext{आणि}
g''(z) & = \Succ(f(z)),
\end{align*}
येथे $f$ वरीलप्रमाणेच आहे: $f(z) = \Succ(\Zero(z))$. आता आपल्याकडे~$h'$
आल्यावर, पुन्हा संयोजन वापरून $h(y) =
h'(\Proj{1}{0}(y),\Proj{1}{0}(y))$ असे ठेवता येते. यावरून मूलभूत
फलनांपासून संयोजने आणि आदिम पुनरावर्तने यांचा अनुक्रम वापरून $h$ ची
व्याख्या करता येते; म्हणून $h$ हे आदिम पुनरावर्ती आहे.
\end{ex}

\end{document}
